% MRM structured abstract: Purpose / Methods / Results / Conclusion. <=250 words.
% Self-contained: no citations, no equation references.
\begin{abstract}
\noindent
\textbf{Purpose:} Surface relaxivity and time-dependent diffusion are two readouts
of the same wall collisions on one substrate. The transverse rate microstructure imaging
fits is a bulk rate plus a surface rate $\rho\,(S/V)$; intra- and extra-axonal water carry
\emph{different} $S/V$, so their $T_2$ differ, biasing any compartment estimate that
normalises by a TE-weighted $b=0$. We quantify the resulting bias on the diffusion
intra-axonal signal fraction and on the myelin water fraction (MWF).

\textbf{Methods:} Closed forms for the interior (Brownstein--Tarr) and exterior
(Novikov--Burcaw) surface rates over myelinated cylinders are derived and validated by
wall-counting Monte Carlo. A calibre-distributed forward model, and grid-free closed forms,
give the intra-axonal fraction ($f_{\mathrm{intra}}$) bias vs packing and echo time and the
MWF bias vs calibre; a within-subject demyelination trajectory traces the longitudinal case.

\textbf{Results:} The interior/exterior $S/V$ ratio is $(1-\mathrm{VF})/(g\,\mathrm{VF})$,
independent of the calibre distribution, crossing unity at fibre volume fraction
$\mathrm{VF}^{\ast}=1/(1+g)$. Physiological white matter sits above it, so surface relaxivity
\emph{over}-weights the intra-axonal signal --- the leading-order intra-axonal fraction bias ---
by ${\approx}12\%$ over the robust packing band at clinical PGSE ($\mathrm{TE}=80$\,ms,
cited $\rho$). That fractions are TE-dependent is known; what is new is the surface-relaxivity
and packing \emph{attribution} and the closed-form sign law, with a testable
packing-dependent TE drift. The same physics reads through relaxometry as a smaller MWF bias:
the thinnest axons' water crosses below the myelin window and is counted as myelin, so fine WM
reads myelin-richer (${\sim}0.33$\,pp at cited $\rho$, far beneath single-voxel noise and
exposed only by regional averaging, but super-linear in the order-of-magnitude-uncertain $\rho$). In primary (lumen-preserving)
demyelination the MWF bias is lumen-set, nearly constant, and cancels in the change.

\textbf{Conclusion:} One $S/V$ surface-relaxivity physics biases both diffusion and
relaxometry microstructure estimates; the $f_{\mathrm{intra}}$ bias is first-order and
packing-dependent (a spatially structured systematic, scaling with the uncertain $\rho$),
the MWF bias small and structural.

\vspace{1em}
\noindent\textbf{Keywords:} intra-axonal signal fraction; myelin water fraction; surface
relaxivity; transverse relaxation; Monte Carlo simulation; white-matter microstructure
\end{abstract}
